% ===================================================================
%  TAULUKKO N:o 16 — Suuri tavaratalo. — Basaari. — Lelut.
%  Traduction 2026 d'apres texto/io, controlee sur texto/fr et
%  texto/sv. Consigne : voir texto/fi/10-taulukko-01.tex.
%
%  Le tableau d'astronomie (bloc 15-1) porte des renvois a LETTRES,
%  (a) a (f), graves sur le tableau lui-meme. L'ido saute la lettre
%  (d), que le fac-simile grave pourtant : la colonne la rend, comme
%  toutes les autres traductions. Ecart declare dans outils/renvoji.py.
%
%  DERNIER TABLEAU DE LA COLONNE FINNOISE. Bilan des seize : treize
%  inversions, toutes de la meme forme — le genitif prepose la ou le
%  possede porte un numero — et toutes defaites par la meme voie, la
%  relative. Le finnois n'a pas de preposition a substituer : il a
%  absorbe ses prepositions dans ses cas, et il faut donc refaire la
%  phrase plutot que d'echanger un mot.
% ===================================================================

%%K t16-apar-1 apar
\VUsaut{4.0mm}
\VUtitre{15.0pt}{60}{TAULUKKO N:o 16}
\VUsaut{2.0mm}
\VUpk{13.2pt}{40}{Suuri tavaratalo. --- Basaari.}
\VUpk{13.2pt}{40}{Lelut.}
\VUsaut{3.4mm}

%%K t16-01-1 p
1. --- Uusi vuosi lähestyy. Suuret tavaratalot ovat valmistaneet \VUgras{lelujen}
\textsuperscript{(1)} ja uudenvuodenlahjojen näyttelynsä.

%%K t16-01-2 p
Kysyin isoisältäni, veisikö hän minut basaariin katsomaan tuota kaunista
näyttelyä, ja hän suostui siihen.

%%K t16-02-1 p
2. --- Ensin pysähdyimme kauniiden \VUgras{vaakunakilpien} eteen, joissa oli
\VUgras{kivääri}, \VUgras{lippalakki}, \VUgras{sapeli}, \VUgras{miekkavyö} ja
pari \VUgras{olkalaattoja}, tai \VUgras{kypärä}, \VUgras{rintahaarniska}
\textsuperscript{(3)} ja \VUgras{pistooli} \textsuperscript{(4)}. Kaikki se
kiinnosti minua paljon enemmän kuin \VUgras{taikalyhdyt} \textsuperscript{(2)}.

%%K t16-tit-1 sub
\VUsaut{3.0mm}
\VUpk{10.2pt}{40}{Kauniita näkyjä.}
\VUsaut{2.6mm}

%%K t16-03-1 p
3. --- Samassa osastossa seisoi \VUgras{vahtisotilas} \textsuperscript{(5)}
\VUgras{vahtikojussaan}; hän näytti vartioivan kokonaista pahvista
eläintarhaa. Kuinka monta eläintä siellä olikaan : \VUgras{papukaija}
\textsuperscript{(6)} orrellaan, \VUgras{sarvikuono} \textsuperscript{(7)},
jolla oli sarvi nenässä, \VUgras{poro} \textsuperscript{(8)},
\VUgras{strutsi} \textsuperscript{(9)}, \VUgras{apina} \textsuperscript{(10)},
joka rikkoi pähkinää, \VUgras{kettu} \textsuperscript{(11)}, joka kantoi pois
kanaa, \VUgras{krokotiili} \textsuperscript{(12)}, joka ammotti valtavine
kitoineen, \VUgras{kameli} \textsuperscript{(13)}, sulava
\VUgras{vinttikoira} \textsuperscript{(14)}, \VUgras{pantteri}
\textsuperscript{(15)}, sitten kaksi eläintä, joita en tuntenut ja jotka,
sanotaan, ovat \VUgras{lumikko} \textsuperscript{(16)} ja \VUgras{hyeena}
\textsuperscript{(17)}. Siellä oli myös \VUgras{leopardi}
\textsuperscript{(18)} ja \VUgras{susi} \textsuperscript{(21)}; aivan lähellä
oli näppäriä \VUgras{rottia} \textsuperscript{(19)} ja pikkuruisia
\VUgras{hiiriä} \textsuperscript{(20)} kumista.

%%K t16-03-1 p suite
Lopuksi \VUgras{virtahepo} \textsuperscript{(22)} ja kyttyräselkäinen
\VUgras{dromedaari} \textsuperscript{(23)} päättivät kokoelman. Olisin jäänyt
sinne vielä useiksi tunneiksi, ellei aivan lähellä olisi ollut komeaa leiriä
täynnä \VUgras{tinasotilaita} \textsuperscript{(29)}, joka veti puoleensa
huomioni.

%%K t16-tit-2 sub
\VUsaut{4.0mm}
\VUpk{10.2pt}{40}{Sotilaat ja sota.}
\VUsaut{2.8mm}

%%K t16-04-1 p
4. --- Isoisäni, joka on \VUgras{invalidi} \textsuperscript{(24)} ja jonka täytyi
jättää armeija \VUgras{haavan} vuoksi, josta hän sai \VUgras{sotamitalin},
kertoo mielellään minulle niistä \VUgras{sotaretkistä}, joihin hän osallistui.
Hän oli onnellinen selittäessään pienen armeijan eri liikkeet pienoiskoossa.
Ryhmä oikeita sotilaita, jotka kävivät basaarissa : \VUgras{pioneeri}
\textsuperscript{(25)}, \VUgras{alppijääkäri} \textsuperscript{(26)}, jonka pää
oli peitetty \VUgras{baskerilla}, jalkaväen \VUgras{vääpeli}
\textsuperscript{(27)} ja \VUgras{tykistön kersantti} \textsuperscript{(28)},
jolla oli kultainen \VUgras{kaluuna} hihassaan, pysähtyi viereemme kuuntelemaan
selityksiä.

%%K t16-05-1 p
5. --- Isoisäni, hyvin ylpeänä siitä, että hänellä oli niin loistava
kuulijakunta, sepitti hetkessä täydellisen taistelusuunnitelman.

%%K t16-05-2 p
Linnoitettua kaupunkia \VUgras{valleineen} ja \VUgras{vallihautoineen} piiritti
\VUgras{vihollisarmeija}, joka pitkän \VUgras{pommituksen} jälkeen oli valmis
\VUgras{rynnäkköön}. Mutta \VUgras{piiritetyt}, joita nälkä pakotti, olivat
uskaltautuneet \VUgras{hyökkäysretkelle} \VUgras{piirittäjien}
\VUgras{leiriin}. Torjuttuaan \VUgras{hyökkäyksen} itsepintaisella
\VUgras{taistelulla} \VUgras{telttaleirin} ympärillä, \VUgras{voittavat}
piirittäjät lähettivät \VUgras{eskadroonansa} ajamaan takaa
\VUgras{voitettuja} hyökkääjiä. Nämä \VUgras{vetäytyivät} ja peittivät
\VUgras{taistelukentän} \VUgras{kuolleilla} ja \VUgras{haavoittuneilla}.
\VUgras{Sairaankantajat} veivät viimeksi mainitut \VUgras{ambulanssiin}
antaakseen \VUgras{kirurgin} hoitaa heitä.

%%K t16-05-3 p
Muutamien \VUgras{pataljoonien} sankarillisesta vastarinnasta huolimatta, jotka
olivat muodostaneet \VUgras{karreen}, \VUgras{tappio} oli täydellinen, armeijan
loput \VUgras{pakenivat} ja jättivät jälkeensä monta \VUgras{vankia}.
Vihollinen meni täydentämään \VUgras{voittoaan} miehittämällä kaupungin.

%%K t16-05-3 p suite
Ehkä siitä tulee sotaretken loppu, ja ehkä on \VUgras{aselevon} jälkeen
mahdollista allekirjoittaa \VUgras{rauhansopimus}.

%%K t16-tit-3 sub
\VUsaut{4.4mm}
\VUpk{10.2pt}{40}{Lahjoja uudeksi vuodeksi.}
\VUsaut{2.8mm}

%%K t16-06-1 p
6. --- Nuoret sotilaat, jotka nauroivat kuunnellessaan veteraanin
maalauksellista kertomusta, pyysivät häneltä vielä muutamia selityksiä. Mitä
minuun tulee, kiertelin puodissa katsellakseni kaunista \VUgras{varsijousta}
\textsuperscript{(32)} ja \VUgras{jousta} \textsuperscript{(33)}
\VUgras{nuolineen} \textsuperscript{(31)}. Sieltä löysin toisen kokoelman
eläimiä : kaksi \VUgras{laululintua} \textsuperscript{(34)} : automaattisen
\VUgras{satakielen} ja \VUgras{kertun}, joka lauloi täydestä kurkusta, ja
sarjan outoja eläimiä : \VUgras{lepakon} \textsuperscript{(35)},
\VUgras{boakäärmeen} \textsuperscript{(36)}, \VUgras{kilpikonnan}
\textsuperscript{(37)}, \VUgras{hylkeen} \textsuperscript{(38)},
\VUgras{valaan} \textsuperscript{(39)} ja \VUgras{hain}
\textsuperscript{(40)} kullatusta nahasta.

%%K t16-07-1 p
7. --- En pysähtynyt kauaksi aikaa katselemaan niitä, sillä aloin nähdä sen,
mistä uneksin : komean \VUgras{nukketeatterin} \textsuperscript{(41)} loistavine
\VUgras{lavasteineen} ja ihastuttavine punaisine \VUgras{esirippuineen}.
\VUgras{Näyttämöllä} kaksi näyttelijää näytti esittävän \VUgras{osaa} hyvin
kiehtovassa \VUgras{näytelmässä}, sillä pienet pahviset \VUgras{katsojat},
jotka olivat sijoitettuina aitioihin tai istuivat \VUgras{nojatuoleilla} ja
\VUgras{permannolla} \VUgras{kapellimestarin} takana, taputtivat vilkkaasti.

%%K t16-07-2 p
Ikävä kyllä tuo teatteri maksoi hyvin paljon.

%%K t16-08-1 p
8. --- Muutoinkin oli vaikea valita leluja, sillä näyteikkunoihin oli kasattu
kaikenlaisia leluja : \VUgras{Harlekiini} \textsuperscript{(42)},
\VUgras{Pulcinella} \textsuperscript{(43)}, \VUgras{Pierrot}
\textsuperscript{(44)}, \VUgras{pöllöjä} \textsuperscript{(45)}, jotka
läiskyttivät siipiään, viheltäviä \VUgras{mustarastaita}
\textsuperscript{(46)}, haukkuvia \VUgras{villakoiria}, \VUgras{fonografi}
\textsuperscript{(47)} ja \VUgras{kärsivällisyyspelejä}.

%%K t16-noto-1 noto
\VUnotes{143.2mm}{%
\textit{(*) Katso taulukko n:o 6.}%
}

%%K t16-08-1 p suite
Yksi noista leluista huvitti minua kovasti : se esitti \VUgras{meritaistelua}
\textsuperscript{(48)}, jossa \VUgras{merirosvot} hyökkäsivät \VUgras{laivan}
kimppuun \VUgras{kookospalmuja} kasvavan maan luona.

%%K t16-09-1 p
9. --- Saatoin hyvin helposti katsella kaikkia noita ihmeitä, sillä tavaratalossa
oli vain vähän ihmisiä, tuskin muutamia kuljeskelijoita, \VUgras{rakuuna}
\textsuperscript{(49)} ja \VUgras{perimismies} \textsuperscript{(50)}, joka
esitti vekselin pää\VUgras{kassalla} \textsuperscript{(51)}.

%%K t16-10-1 p
10. --- Kysyin isoisältäni, mihin käytetään niitä kirjoja, jotka ovat hyllyillä
\VUgras{kassanhoitajattaren} takana; hän vastasi minulle, että ne ovat
\VUgras{tilikirjoja}.

%%K t16-tit-4 sub
\VUsaut{3.6mm}
\VUpk{10.2pt}{40}{Puodin edessä töllistelyä.}
\VUsaut{2.8mm}

%%K t16-11-1 p
11. --- Kun jätimme leluosaston, menimme satulasepänliikkeeseen luodaksemme
silmäyksen \VUgras{satuloihin} \textsuperscript{(53)}, \VUgras{jalustimiin}
\textsuperscript{(54)}, \VUgras{ohjaksiin} \textsuperscript{(55)},
\VUgras{kuolaimiin} \textsuperscript{(56)} ja \VUgras{ratsupiiskoihin}.
Samalla kun \VUgras{osastonjohtaja} \textsuperscript{(57)} antoi isoisälleni
muutamia tietoja, menin katselemaan \VUgras{posliinin} ja \VUgras{kristallin}
\textsuperscript{(58)} näyttelyä, jossa on kauniita \VUgras{salaattikulhoja} ja
hienoja \VUgras{teekannuja} \textsuperscript{(59)}.

%%K t16-12-1 p
12. --- Mennäksemme ensimmäiseen kerrokseen kuljimme \VUgras{korsettien}
\textsuperscript{(60)} näyteikkunan takaa, sukkaosaston vierestä, jossa hyvin
kauniita \VUgras{alushousuja} \textsuperscript{(63)} oli näytteillä. Aivan
lähellä \VUgras{myyjätär} \textsuperscript{(61)} antoi \VUgras{ostajattaren}
\textsuperscript{(62)} valita kauniita silkkisiä \VUgras{kankaita}
\textsuperscript{(64)}.

%%K t16-12-2 p
Noustessani portaita huomasin, että tavaratalo on koristettu japanilaisin
\VUgras{lyhdyin} \textsuperscript{(65)}, \VUgras{päivänvarjoin}
\textsuperscript{(66)} ja valtavin \VUgras{palmuin} \textsuperscript{(70)}.

%%K t16-13-1 p
13. --- Ensimmäisessä kerroksessa ei ollut paljon nähtävää; vain huonekaluja (*),
\VUgras{kassakaappeja} \textsuperscript{(67)}, \VUgras{lipastoja}
\textsuperscript{(68)}, \VUgras{lastenvaunuja} \textsuperscript{(69)}. Mutta
tavaratalon kokonaisnäkymä oli hyvin \VUgras{huvittava}.

%%K t16-14-1 p
14. --- Jalkojemme juuressa näin soittimet : \VUgras{harppuja}
\textsuperscript{(71)}, \VUgras{kitaroita} \textsuperscript{(72)},
\VUgras{mandoliineja} \textsuperscript{(73)}, \VUgras{venttiilikorneteja}
\textsuperscript{(74)}, \VUgras{klarinetteja} \textsuperscript{(75)}, viuluja
\VUgras{jousineen} \textsuperscript{(76)} ja jopa \VUgras{ison rummun}
\textsuperscript{(77)}. Hieman kauempana, paperiosaston luona,
\VUgras{ylioppilas} \textsuperscript{(78)} tinki \VUgras{myyjältä}
\textsuperscript{(79)} vihkokansiota. Heidän vieressään erotin kauniin
\VUgras{aapisen} \textsuperscript{(80)}.

%%K t16-14-2 p
Keskellä tavarataloa oli pystytetty korkea \VUgras{joulukuusi}
\textsuperscript{(81)}, joka oli kokonaan varustettu kynttilöillä, koristeilla ja
kaikenlaisilla leluilla : \VUgras{puuhevosilla} \textsuperscript{(82)},
\VUgras{nukeilla} \textsuperscript{(83)}, j. n. e..

%%K t16-14-3 p
\VUgras{Hajuvesiosasto} \textsuperscript{(84)} \VUgras{hammasaineineen} ja
erilaisine \VUgras{hajuvesineen} levitti hyvin hyvää tuoksua, joka tuli meitä
kohti.

%%K t16-tit-5 sub
\VUsaut{3.4mm}
\VUpk{10.2pt}{40}{Ostamme jotakin.}
\VUsaut{2.8mm}

%%K t16-15-1 p
15. --- Koska isoisäni alkoi tuntea ajan pitkäksi, laskeuduimme suuria portaita
alas. Hän pysähtyi hetkeksi \VUgras{tähtitieteellisen taulukon}
\textsuperscript{(85)} eteen, joka, hän sanoi minulle, esittää
\VUgras{komeettoja} \textsuperscript{(a)}, \VUgras{kuun- ja
auringonpimennyksiä} \textsuperscript{(b)}, tuuliruusun \VUgras{neljine
ilmansuuntineen} \textsuperscript{(c)} \VUgras{(Pohjoinen, Etelä, Itä ja
Länsi)}, kartan kahdesta \VUgras{pallonpuoliskosta} \textsuperscript{(d)},
\VUgras{planeetat} \textsuperscript{(e)} ja \VUgras{aurinkokunnan}
\textsuperscript{(f)}. En ymmärtänyt mitään kaikesta siitä, ja minua kiehtoi
paljon enemmän kaluunoilla koristettu livree, joka oli
\VUgras{juoksupojalla} \textsuperscript{(86)}, joka kulki ohi, ja kauniit
\VUgras{viuhkat} \textsuperscript{(87)}, jotka olivat vieressäni,
\VUgras{kukkaroiden} \textsuperscript{(88)} ja \VUgras{lompakoiden}
\textsuperscript{(89)} luona.

%%K t16-16-1 p
16. --- Ihailin myös näyttelyssä \VUgras{höyheniä} \textsuperscript{(90)} ja
\VUgras{vyönsolkia} \textsuperscript{(91)}, kauniita \VUgras{tekokukkia}
\textsuperscript{(94)}, jotka oli sulavasti aseteltu koreihin.

%%K t16-16-1 p suite
\VUgras{Orvokit}, \VUgras{leukoijat}, \VUgras{hyasintit}
\textsuperscript{(95)}, \VUgras{pelargonit} \textsuperscript{(96)},
\VUgras{liljat} \textsuperscript{(97)} ja \VUgras{krassit}
\textsuperscript{(98)} olivat hyvin hyvin jäljiteltyjä.

%%K t16-tit-6 sub
\VUsaut{4.4mm}
\VUpk{10.2pt}{40}{Isoisäni lahja.}
\VUsaut{2.8mm}

%%K t16-17-1 p
17. --- Pyysin isoisääni ostamaan minulle yhden niistä \VUgras{hammasharjoista}
\textsuperscript{(92)}, jotka olivat rasiassa \VUgras{hiusneulojen}
\textsuperscript{(93)} vieressä. Hän suostui, ja menin itse maksamaan ostokseni
kassalle, jonka vieressä valmistettiin tärkeää \VUgras{lähetystä}
\textsuperscript{(100)} eräälle \VUgras{asiakkaalle} \textsuperscript{(99)}.

%%K t16-18-1 p
18. --- Äkkiä syntyi lievä hämminki niiden ihmisten keskuudessa, jotka olivat
pysähtyneet taiteellisten \VUgras{pronssitöiden} \textsuperscript{(101)} luo.
\VUgras{Varas} \textsuperscript{(102)} oli juuri otettu kiinni itse teosta
\VUgras{tarkastajan} \textsuperscript{(103)} toimesta, kun hän koetti varastaa
joltakulta. Huolehdittiin siitä, että haettiin poliisikonstaapeli
\VUgras{pidättämään} hänet, ja juuri kun menimme ulos, tekijä vietiin
vankilaan.

\VUsaut{6.0mm}
\VUfilet{22mm}
